% Autorska priprava k Realisticky.cz 020313–020314; ulohy jsou autorske a obmenene.
% Zamerne nenahrazuje vyklad presnymi kopiemi cizich textu/prikladu.
\documentclass[11pt,a4paper]{article}
\input{../../sablony/pripravy/priprava.sty}
\usepackage{tikz}
\nastavPripravu{1. ročník}{Rovnice a nerovnice}{85}{Soustavy $n$ rovnic o $n$ neznámých: princip}
\begin{document}
\blok{Návaznost a motivace}
U dvou neznámých jsme snižovali počet neznámých v rovnici. Stejný nápad funguje i u tří, čtyř nebo obecně $n$ neznámých: postupnou eliminací získáme jednodušší soustavu a potom dosazujeme zpět.
\blok{Princip a odvození}
Pro soustavu $n\times n$ používáme stejnou myšlenku jako u sčítací metody. Cílem je postupně odstranit první neznámou z dalších řádků, druhou z následujících atd. Přítomnost $n$ rovnic však \emph{nezaručuje} právě jedno řešení: rovnice mohou být závislé či ve sporu. Maticový formalismus je pouze poznámka navíc.

\begin{center}\small
\(\underbrace{3\ \text{rovnice},\ 3\ \text{neznámé}}_{\text{eliminace}}\longrightarrow
\underbrace{2\ \text{rovnice},\ 2\ \text{neznámé}}_{\text{eliminace}}\longrightarrow
1\ \text{rovnice},\ 1\ \text{neznámá}\)
\end{center}
\blok{Řešené příklady: od jednoduchého ke složitějšímu}
\begin{prikladReseniBox}[title={Příklad 1}]
Vyřešte eliminačně $\begin{cases}2x+y-z=4,\\x-y+z=2,\\x+y+z=6.\end{cases}$
\tcblower
\begin{align*}
(2x+y-z)+(x-y+z)&=4+2 \\
 3x&=6\Rightarrow x=2 \\
 -y+z&=0 \\
 y+z&=4 \\
 y&=2,\quad z=2
\end{align*}

Zkouška všech tří rovnic dává postupně $4,2,6$.
\end{prikladReseniBox}
\begin{prikladReseniBox}[title={Příklad 2}]
Vyřešte trojúhelníkovou soustavu $\begin{cases}x+y+z=6,\\2x+y+z=8,\\x+2y+z=9.\end{cases}$
\tcblower
\begin{align*}
(2x+y+z)-(x+y+z)&=8-6 \\
 x&=2 \\
 (x+2y+z)-(x+y+z)&=9-6 \\
 y&=3,\quad z=1
\end{align*}
\end{prikladReseniBox}
\poznamka{Smyslem není memorovat tabulku elementárních řádkových operací. Stačí rozumět eliminaci a zpětnému dosazení.}
\blok{Samostatné procvičování}
\begin{zadaniReseniBox}[title={Úloha 1}]
$\begin{cases}x+y+z=9,\\y+z=5,\\z=2.\end{cases}$
\tcblower
\begin{align*}
z&=2 \\
 y&=5-2=3 \\
 x&=9-3-2=4
\end{align*}
\end{zadaniReseniBox}
\begin{zadaniReseniBox}[title={Úloha 2}]
Co se změní, když máme jen $x+y=4$, $y+z=5$?
\tcblower
\begin{align*}
y&=t \\
 x&=4-t \\
 z&=5-t,\quad t\in\mathbb R
\end{align*}

Bez další nezávislé podmínky je nekonečně mnoho řešení.
\end{zadaniReseniBox}
\begin{zadaniReseniBox}[title={Úloha 3}]
Vyřešte $\begin{cases}x+y+z=7,\\x+y=4,\\y+z=5.\end{cases}$
\tcblower
\begin{align*}
z&=7-4=3 \\
 y&=5-3=2 \\
 x&=4-2=2
\end{align*}
\end{zadaniReseniBox}
\blok{Volitelná rezerva}
\begin{zadaniReseniBox}[title={Úloha 1}]
Mohou mít tři rovnice se třemi neznámými nulový počet řešení?
\tcblower
\begin{align*}
x+y+z&=5 \\
 x+y+z&=6 \\
 5&=6
\end{align*}

Ano. Dvě podmínky si odporují a třetí rovnice to neodstraní.
\end{zadaniReseniBox}
\blok{Domácí úkol}
\begin{zadaniReseniBox}[title={Úloha 1}]
$\begin{cases}x+y+z=8,\\x+y=5,\\y+z=4.\end{cases}$
\tcblower
\begin{align*}
z&=8-5=3 \\
 y&=4-3=1 \\
 x&=5-1=4
\end{align*}
\end{zadaniReseniBox}
\begin{zadaniReseniBox}[title={Úloha 2}]
$\begin{cases}x+y+z=6,\\x+y=3,\\y+z=4.\end{cases}$
\tcblower
\begin{align*}
z&=6-3=3 \\
 y&=4-3=1 \\
 x&=3-1=2
\end{align*}
\end{zadaniReseniBox}
\begin{zadaniReseniBox}[title={Úloha 3}]
$\begin{cases}x+y+z=5,\\x+y=2,\\y+z=3.\end{cases}$
\tcblower
\begin{align*}
z&=5-2=3 \\
 y&=3-3=0 \\
 x&=2-0=2
\end{align*}
\end{zadaniReseniBox}
\poznamka{Orientační časové rozvržení: 5 + 12 + 12 + 10 + 6 = 45 min. Přednost má pochopení jádra; rezerva není součástí časového součtu. Domácí úlohy jsou určeny na samostatnou práci, řešení zůstávají učiteli.}
\end{document}
